\section{The successive viable-set construction}
\label{sec:lvv_successive_construction}

Fix a local dimension $q\ge2$ and a gap parameter $\Delta>0$.
The construction below follows the induction of \cite{Landau2015Polynomial},
with the independent trimming radius corrected as
in~\cite{gap:lvv13_trimming_radius}. The cardinality constants use the
stated constant-cut estimate and corrected net exponent
of~\cite{gap:lvv13_constant_bond}. Error reduction uses the finite Gaussian
construction of \cite[Section~5.4, Lemmas~18--19]{Huang2015GroundState}.
This permits arbitrary ground energy, unlike the uniform polynomial
estimate for the original power sampler; the distinction is explained
in~\cite{gap:lvv13_sampling_ground_energy}.
All polynomial bounds have the quantitative scope specified in those constructions.

\begin{definition}[Successive size and bond budgets]
    \label{def:lvv_successive_budgets}
    \lean{TNLean.MPS.AreaLaw.lvvSuccessiveError,
        TNLean.MPS.AreaLaw.lvvSizeTrimmingDegree,
        TNLean.MPS.AreaLaw.lvvSizeTrimmingCoefficient,
        TNLean.MPS.AreaLaw.lvvSuccessiveCardBound,
        TNLean.MPS.AreaLaw.lvvSuccessiveBondBound,
        TNLean.MPS.AreaLaw.lvvSuccessiveError_pos}
    \leanok
    \uses{thm:lvv_size_bond_trimmed_family, def:huang_filter_parameters}
    Put $c_\Delta=(\Delta/169)^2>0$. Let $D$ be the constant-cut bond bound
    used in size trimming at precision $c_\Delta$, and define
    \begin{align}
        p   & =2(qD)^2,                                           &
        A   & =\left\lceil\frac{2(qD)^2}{c_\Delta}\right\rceil_+,   \\
        C_s & =(D+1)D(2A+1)^p,                                    &
        m_L & =(48L)^2r_{\rm ref}(L).
    \end{align}
    Here $r_{\rm ref}(L)$ is the reference bond from the existing constant-accuracy
    MPS approximation at length $48L$. Write $E_L$ and $R_L$ for the energy count
    and finite-filter rank budget with precision parameter $c_\Delta$. Set
    \begin{align}
        s_L=C_sL^pE_LR_L,\qquad b_L=R_L\max\{1,m_L\}.
    \end{align}
\end{definition}

\begin{lemma}[Uniform polynomial budgets]
    \label{thm:lvv_successive_polynomial_budgets}
    \lean{TNLean.MPS.AreaLaw.exists_lvvSuccessive_uniform_polynomial}
    \leanok
    \uses{def:lvv_successive_budgets, thm:huang_filter_uniform_polynomial,
        lem:lvv_bond_trimming_radius_polynomial}
    There are natural numbers $P,d$, with $P\ge1$, depending only on $q,\Delta$
    and chosen before the length, such that for every $L\ge1$,
    \begin{align}
        \max\{1,s_L\}    & \le P(L+1)^d, &
        \max\{1,b_L\}    & \le P(L+1)^d, &
        s_L\max\{1,b_L\} & \le P(L+1)^d.
    \end{align}
\end{lemma}
\begin{proof}
    \leanok
    The energy count and rank product have a uniform natural-coefficient
    polynomial bound. So does the corrected reference trimming radius, after
    taking a natural ceiling of its coefficient. Multiply these bounds by the
    fixed size coefficient and the power $L^p$. Products of the resulting two
    polynomial bounds control both budgets and their final-span product. Taking
    maxima with one also covers the scalar initial family.
\end{proof}

\begin{lemma}[A complete successive step]
    \label{thm:lvv_complete_successive_step}
    \lean{TNLean.MPS.AreaLaw.exists_positive_chain_successive_step}
    \leanok
    \uses{def:lvv_successive_budgets, thm:lvv_size_bond_trimmed_family,
        thm:huang_positive_chain_operator, thm:huang_positive_chain_viable_filter}
    Let $L\ge1$ and $1\le i\le L$. Let
    $H=\sum_{j=0}^{L-2}h_j$, where $h_j$ is supported on sites $j,j+1$ and
    $0\le h_j\le I$. Given a finite family $S$ on the first $i-1$ sites whose
    members have a common finite MPS bond bound, one can choose a family $T$
    on the first $i$ sites with $|T|\le s_L$ and individual effective bond
    bound $b_L$, before any ground energy or ground vector is supplied.
    For every unit ground vector $\Gamma$ with projector gap $\Delta$,
    $(i-1,c_\Delta/L)$-viability of $S$ implies $(i,c_\Delta/L)$-viability of $T$.
\end{lemma}
\begin{proof}
    \leanok
    Extend each member by the one-site basis. Apply the actual size-trimming
    construction to this extended family, using the exact adjacent-pair
    embedding of $H$. Independently trim every retained vector to bond $m_L$.
    The resulting family has cardinality at most $C_sL^p$ and is half-viable.
    Apply the finite Gaussian filter family. Its cardinality multiplier is
    $E_LR_L$, its individual bond multiplier is $R_L$, and it restores
    viability precision $c_\Delta/L$. Each family is selected before the
    ground pair, so the composition has the same quantifier order.
\end{proof}

\begin{lemma}[One recursively chosen viable sequence]
    \label{thm:lvv_coherent_successive_sequence}
    \lean{TNLean.MPS.AreaLaw.exists_positive_chain_successive_viable_sets}
    \leanok
    \uses{thm:lvv_complete_successive_step}
    Under the physical hypotheses of the preceding theorem, there is one
    sequence $S_0,\ldots,S_L$, chosen before every ground pair, with $S_0$ the
    scalar singleton. For each positive cut, $|S_i|\le s_L$ and every member
    has effective MPS bond bound $b_L$. For every unit ground vector $\Gamma$
    with projector gap $\Delta$, all these same families are
    $(i,c_\Delta/L)$-viable.
\end{lemma}
\begin{proof}
    \leanok
    Carry a family together with its cardinality, individual bond and all-ground
    viability assertions. At zero take the scalar singleton, whose viable
    witness is $\Gamma$ itself. Define the successor by choosing the complete
    step applied to the actual preceding family. Its viability assertion follows
    from that of its predecessor. Natural-number recursion therefore gives a
    single successive run. No ground pair enters any successor choice.
\end{proof}

\begin{theorem}[Uniform polynomial successive families and final-span MPS]
    \label{thm:lvv_uniform_successive_approximation}
    \lean{TNLean.MPS.AreaLaw.exists_uniform_positive_chain_successive_approximation}
    \leanok
    \uses{thm:lvv_successive_polynomial_budgets, thm:lvv_coherent_successive_sequence,
        thm:lvv_left_span_bond}
    Fix $q\ge2$ and $\Delta>0$. There exist natural numbers $P,d$, with $P\ge1$,
    chosen before $L$ and the interaction, such that the following holds.
    For every $L\ge1$ and every adjacent-pair positive-contraction interaction,
    one can choose a single successive sequence $S_0,\ldots,S_L$ with scalar
    initial family, cardinalities and individual effective bond bounds at most
    $P(L+1)^d$, before every ground pair.

    Suppose $\|\Gamma\|=1$, $H\Gamma=E_0\Gamma$ and
    $H-E_0I\ge\Delta(I-|\Gamma\rangle\langle\Gamma|)$. Then every $S_i$ is
    $(i,c_\Delta/L)$-viable, and there is a vector $\psi$ in the span of the
    actual final family, read on the full chain, such that
    \begin{align}
        \|\psi\|=1,\qquad
        |\langle\psi,\Gamma\rangle|\ge1-c_\Delta/L.
    \end{align}
    This $\psi$ is an open-boundary MPS with bond dimension at most $P(L+1)^d$.
\end{theorem}
\begin{proof}
    \leanok
    Choose the polynomial budgets first and then the recursively constructed
    sequence. Enlarge the explicit budgets to this polynomial at each cut.
    At the final cut, the right configuration space is one-dimensional. Thus
    a normalized viable witness belongs to the span of the actual $S_L$ in
    full-left coordinates. Pull it back by the coordinate isometry, preserving
    norm and overlap. The finite-span rank estimate bounds its MPS bond by
    $s_L\max\{1,b_L\}$, already controlled by the same polynomial.
\end{proof}

The conclusion states the overlap precision $1-c_\Delta/L$ and the displayed
polynomial existence bounds. It makes no computational running-time assertion.
